\section{What a change of formulation can preserve}
\label{sec:formulation}

\Cref{thm:completion} concerns one formulation and one barrier. This
section gives two movement lower bounds for other formulations,
including those with auxiliary variables. For log-determinant barriers
of symmetric cones, the rank of an exposing dual solution bounds primal
movement; for every logarithmically homogeneous barrier, a bound on the
factor dimensions bounds full primal--dual movement. Neither statement derives a primal
lower bound from an upper bound on the barrier parameter.

\subsection{Exposing dual solutions and free auxiliary variables}
Let $K=\prod_iK_i$ be a product of symmetric cones in Euclidean Jordan
algebras with trace inner products, and let
$\Omega=\{x\in\interior K:Ax=b\}\ne\varnothing$.
Suppose the objective $\ell$ has finite supremum $\ell_*$ over
$\Omega$ and an attained dual solution $s=(s_i)\in K$ that exposes
the optimum:
\begin{equation}
 \ell_*-\ell(x)=\sum_i\ip{s_i}{x_i}\qquad(x\in\Omega).
 \label{eq:exposing-certificate}
\end{equation}
Let $c_i$ be the support idempotent of $s_i$, let
$q_i=\rank(s_i)$, and omit factors of rank zero.
Here $P(c_i)$ is the orthogonal Peirce projection onto the algebra with
unit $c_i$, and $\det_{c_i}$ is its determinant.

\begin{theorem}[Distance bound from the rank of an exposing dual solution]
\label{thm:exposed-rank}
Use the barrier
\[ F(x)=-\sum_i\alpha_i\log\det_i x_i,\qquad \alpha_i\geq1, \]
restricted to $\Omega$. Put $Q_\alpha=\sum_i\alpha_iq_i>0$ and, for
any reference $x^c\in\Omega$, define
\begin{equation}
 \Delta_{c,\alpha}=Q_\alpha
 \left[\prod_i
 \frac{\{\det_{c_i}(s_i)\det_{c_i}(P(c_i)x_i^c)\}^{\alpha_i}}
      {\alpha_i^{\alpha_iq_i}}\right]^{1/Q_\alpha}.
 \label{eq:weighted-exposed-scale}
\end{equation}
For $\eps>0$, every feasible $y$ with objective gap at most $\eps$
satisfies
\begin{equation}
 d_F(x^c,y)\geq\sqrt{Q_\alpha}
            \pos{\log(\Delta_{c,\alpha}/\eps)}.
 \label{eq:exposed-distance}
\end{equation}
The bound does not restrict eigenvalues outside the support of $s$ or
auxiliary variables, and needs no strict complementarity. If a method starts at $x_0$ with
$d_F(x_0,x^c)\leq D_0$ and reaches such a $y$ after $N$ rounds,
each moving the iterate a distance at most $B_0>0$, then
\begin{equation}
 N\geq\frac{\pos{\sqrt{Q_\alpha}\log(\Delta_{c,\alpha}/\eps)-D_0}}{B_0}.
 \label{eq:exposed-rounds}
\end{equation}
In particular, if $\norm{x_0-x^c}_{x^c}\leq\rho$ with $0\leq\rho<1$
and each round has at most $m\in\N$ forward $R$-chords, $0<R<1$,
one can take $D_0=-\log(1-\rho)$ and
$B_0=-m\log(1-R)$.
\end{theorem}
\begin{proof}
We give the compression argument in full to fix the metric
normalization, including the exceptional Jordan algebra. For one idempotent $c$ of
rank $q$, put $a=P(c)x$ and $\phi_c(x)=-\log\det_c a$.
The Hessian of $-\log\det x$ is $P(x^{-1})$ with inverse $P(x)$,
where $P(u)=2L_u^2-L_{u^2}$ is the quadratic representation and
$L_uv=u\circ v$. Since $a^{-1}$ belongs to the Peirce algebra,
\[
 d\phi_c(x)[h]=-\ip{a^{-1}}{h},\qquad
 \norm{d\phi_c(x)}_{x,*}^2=\ip{a^{-1}}{P(x)a^{-1}}=q.
\]
Indeed, the fundamental identity $P(P(u)v)=P(u)P(v)P(u)$ with $u=c$
gives $P(c)P(x)P(c)=P(a)$; then $P(a)a^{-1}=a$ and
$\ip{a^{-1}}{a}=q$. The compression $a$ is positive definite in its
Peirce algebra: if $x\succeq\delta e$, positivity of $P(c)$ gives
$a\succeq\delta c$.
For $\Phi=\sum_i\alpha_i\phi_{c_i}$, the product metric therefore gives
$\norm{d\Phi}_{F,*}^2=Q_\alpha$ before affine restriction and at most
$Q_\alpha$ afterwards. These quadratic-representation identities are
standard~\cite{FarautKoranyi1994}; we do not claim the principal-minor
metric identity as new.

Set $g_i=\ip{s_i}{y_i}>0$ on the active factors. Inside the Peirce
algebra, $z_i=P(s_i^{1/2})P(c_i)y_i$ has trace $g_i$ and determinant
$\det_{c_i}(s_i)\det_{c_i}(P(c_i)y_i)$.
Spectral AM--GM gives
\[
 \det_{c_i}(P(c_i)y_i)
 \leq\det_{c_i}(s_i)^{-1}(g_i/q_i)^{q_i}.
\]
Under $\sum_i g_i\leq\eps$, the weighted product
$\prod_i(g_i/q_i)^{\alpha_iq_i}$ is maximized at
$g_i=\eps\alpha_iq_i/Q_\alpha$. Thus
\[
 \prod_i\det_{c_i}(P(c_i)y_i)^{\alpha_i}
 \leq(\eps/Q_\alpha)^{Q_\alpha}
 \prod_i\alpha_i^{\alpha_iq_i}\det_{c_i}(s_i)^{-\alpha_i}.
\]
Comparing with $x^c$ proves
$\Phi(y)-\Phi(x^c)\geq Q_\alpha\log(\Delta_{c,\alpha}/\eps)$.
Integrating this dual-norm bound along any feasible curve proves \eqref{eq:exposed-distance}. The triangle inequality and \cref{lem:chord-distance} prove the round
statements.
\end{proof}
For unit weights this gives $Q=\sum_iq_i$ and
$\Delta_c=Q[\prod_i\det_{c_i}(s_i)\det_{c_i}(P(c_i)x_i^c)]^{1/Q}$.
The coefficient $\sqrt{Q_\alpha}$ is asymptotically sharp on the full
cone. Apply a quadratic-representation
isometry sending the reference point to the unit element, diagonalize
the transformed dual solution, and multiply its support coordinates by
$t$, keeping the others equal to one. The gap is
$t\ip{s}{e}$ and the curve has length
$\sqrt{Q_\alpha}\log(1/t)$. Together with
\eqref{eq:exposed-distance}, this proves the exact leading coefficient
as $\eps\downarrow0$. An affine slice need not contain this path; we
claim no sharp upper bound on an arbitrary slice.

Hauser and G\"uler's classification
\cite[Theorem~5.5]{HauserGuler2002} states that every self-scaled barrier
on a symmetric cone is, up to an additive constant, a sum of these
log-determinants on its irreducible factors with coefficients
$\alpha_i\geq1$. Thus \cref{thm:exposed-rank} covers every such barrier,
with the weighted scale $\Delta_{c,\alpha}$; replacing $Q$ by $Q_\alpha$
but keeping the unweighted $\Delta_c$ would be incorrect.
The bound depends on the rank of $s$; a general lower bound on the
ranks attainable by different lifted formulations is a separate
question, which the theorem neither assumes nor answers.

\subsection{A bound from factor dimensions for arbitrary cone barriers}
\begin{theorem}[Bounded factor dimension]
\label{thm:dimension-dictionary}
Let a compact full-dimensional convex body $C\subset\R^D$, $D\geq1$,
have an exact affine lift $C=\pi(K\cap\mathcal L)$, where $K$ is a
finite product of proper cones, $\mathcal L$ is an affine subspace,
and $\pi$ is an affine map. Let the minimal face be the smallest face
of $K$ containing $K\cap\mathcal L$; it is a product of faces of the
factors of $K$. Regarding this product as a cone in its linear span and
removing its zero factors gives the reduced product cone. Suppose
each of its factors has dimension at most an integer $d\geq2$.
For $h\geq1$ define
\[
 \Psi_d(h)=2\lfloor h/d\rfloor+\min\{h\bmod d,2\}.
\]
Every logarithmically homogeneous standard self-concordant barrier on
the reduced product cone, including one that couples the factors,
satisfies
\begin{equation}
 \nu\geq\Psi_d(D+1)\geq\frac{2(D+1)}d.
 \label{eq:dimension-parameter}
\end{equation}
For a strictly feasible primal--dual conic pair on this cone, consider
routes from an exact central point with gap $\Delta_0$ to a strictly
feasible pair with gap at most $\eps$, $0<\eps<\Delta_0$. If each
round has at most $m\in\N$ forward $R$-chords in the product metric,
$0<R<1$, the number of rounds is at least
\begin{equation}
 \frac{\sqrt{\Psi_d(D+1)/2}\log(\Delta_0/\eps)}{-m\log(1-R)}.
 \label{eq:dimension-movement}
\end{equation}
\end{theorem}
\begin{proof}
Let $M$ be the dimension of the span of the reduced product cone.
The affine projection of its feasible slice spans $\R^D$, so $M\geq D$.
If $M=D$, the linear part of $\pi$ is invertible on this span, so the
feasible slice has full affine hull and the affine equations impose no
restriction: the feasible slice is the whole reduced cone.
Its invertible affine image is unbounded, contrary to compactness.
Therefore $M\geq D+1$. This elementary argument is consistent
with the classical slack-factorization framework~\cite{GouveiaParriloThomas2013},
but needs no choice of boundary certificates.

Let $r_0$ and $q_0$ be the numbers of ray factors and of factors
of dimension at least two. In each nonray factor choose a two-dimensional linear plane
containing an interior point. Its cone section is proper and
full-dimensional in the plane, hence linearly isomorphic to $\R_+^2$;
its relative interior lies in the factor's interior and its boundary in
the factor's boundary. Together with all rays, these sections give an orthant section of
dimension $r_0+2q_0$ through the interior of the reduced cone.
Restriction preserves standard self-concordance, the degree of
logarithmic homogeneity, and boundary blow-up. The classical orthant parameter lower bound
\cite{NesterovNemirovskii1994,Hildebrand2013} therefore gives
$\nu\geq r_0+2q_0$, even when the original barrier couples factors.
The dimension bound $d$ gives $r_0+dq_0\geq M\geq D+1$.

To solve the integer minimization, write $h=ad+b$, $0\leq b<d$.
For $q_0\leq a$ the smallest possible value of $r_0+2q_0$ is
$2q_0+h-dq_0=h-(d-2)q_0$, minimized at $q_0=a$ and equal to $2a+b$.
For $q_0\geq a+1$ this value is at least $2a+2$, attained with no
rays. The minimum is therefore $2a+\min(b,2)$; the case $b=0$ is
attained with $q_0=a$. Since $\min(b,2)\geq2b/d$, this proves
\eqref{eq:dimension-parameter}. Apply \cref{thm:pd-gap-set} for
\eqref{eq:dimension-movement}.
\end{proof}
The parameter bound concerns barriers on the reduced product cone, not
arbitrary barriers on $C$ or on the affine slice. The theorem combines
classical restriction, dimension, and primal--dual arguments; it does not determine the least dimension of a lift. For a
product of balls, $D=\sum_as_a$, so the bound uses the total ball
dimension and no curvature or regularity assumptions.

\subsection{Grouped balls, norm trees, and PSD packing}
\label{subsec:concrete-formulations}
The following formulations show that fewer conic blocks need not
reduce movement in the given barrier metric. Full proofs,
including all auxiliary variables, are in \cref{app:formulations}.

Take $k$ unit balls, a unit vector $c_a$ for each, and the objective
$\ell=\sum_a\ip{c_a}{w_a}$ with maximum $k$.
For a grouped paraboloid lift, partition each ball's coordinates into
$h$ groups, introduce $q_{aG}=s_{aG}-\norm{w_{aG}}^2>0$ with
$\sum_Gs_{aG}=1$, and use $F=-\sum_{a,G}\log q_{aG}$.
For a PSD packing, place the
$H=kh$ group vectors as columns into real matrix blocks $W_\ell$ in any
arrangement, let every off-diagonal entry of $S_\ell$ be free, and use
\[
 D_\ell=S_\ell-W_\ell^TW_\ell\succ0,\qquad
 F=-\sum_\ell\log\det D_\ell,
 \qquad \sum_{\gamma\text{ from ball }a}(S_{\ell(\gamma)})_{\gamma\gamma}=1.
\]
A group column may be padded with zeros; its free entries are its
group's coordinates. Each group occurs exactly once.
Both barriers are standard self-concordant with exact restricted
parameter $H$; for each, the central path
$z(\tau)=\argmin(F-\tau\ell)$ has
$w_a=r(\tau)c_a$, with
\begin{equation}
 q(\tau)=\frac{2}{\sqrt{h^2+\tau^2}+h},\quad
 r(\tau)=\frac{\tau}{\sqrt{h^2+\tau^2}+h},\quad
 \left\|\frac{dz}{d\log\tau}\right\|^2
   =H\left(1-\frac h{\sqrt{h^2+\tau^2}}\right).
 \label{eq:grouped-packed-center}
\end{equation}
On the path, $q_{aG}=q(\tau)$ (grouped lift) and
$D_\ell=q(\tau)I$ (packing). For every $\tau_0\geq0$ and every
feasible $y$ with gap at most $\eps$,
\begin{equation}
 d_F(z(\tau_0),y)\geq\sqrt H\,
          \pos{\log(q(\tau_0)H/(2\eps))}.
 \label{eq:grouped-packed-distance}
\end{equation}
At the analytic center $\tau_0=0$ the argument of the logarithm is
$k/(2\eps)$, independent of the packing and its off-diagonal
variables.

For a rooted norm tree with $b$ internal nodes, each with at least two
children, use one Lorentz-cone block per internal node:
\begin{equation}
 q_{av}=t_{av}^2-\sum_{u\in I(v)}t_{au}^2
                    -\sum_{j\in J(v)}w_{aj}^2>0,
 \qquad t_{av}>0,
 \qquad t_{a,\mathrm{root}}=1.
 \label{eq:norm-tree-domain}
\end{equation}
Here $I(v)$ and $J(v)$ are the internal and leaf children of $v$.
The positive sheet $t_{av}>0$ is needed for convexity. Use $F=-\sum_{a,v}\log q_{av}$ and
put $L=kb$. Its exact parameter is $\nu=k\nu_T$, where
\begin{equation}
 \nu_T=\begin{cases}2b-2,&\text{the root has only internal children},\\
                         2b-1,&\text{the root has a leaf child}.
          \end{cases}
 \label{eq:tree-parameter}
\end{equation}
Its central path is \eqref{eq:grouped-packed-center} with $h=b,H=L$,
regardless of the tree shape. If $n_v$ counts internal
nodes in the subtree at $v$, including $v$, and $C_{av}=\sum_{j\text{ below }v}c_{aj}^2$, the
internal coordinates are $t_{av}^2=n_vq+r^2C_{av}$.
At the analytic center $\bar z$, $q_{av}=1/b$, and every feasible $y$ with gap at most
$\eps$ satisfies the sharper barrier-value bound
\begin{equation}
 d_F(\bar z,y)\geq\frac{L}{\sqrt\nu}
                        \pos{\log(k/(2\eps))}
 \geq\sqrt{L/2}\,\pos{\log(k/(2\eps))}.
 \label{eq:tree-distance}
\end{equation}
\Cref{app:formulations} proves the exact parameter
\eqref{eq:tree-parameter} by a metric calculation; simply subtracting
the fixed root's contribution from the cone parameter would not be
justified.

For all three formulations, the central path from the analytic center
to objective gap $\eps$ has length
\begin{equation}
 \sqrt{M_0}\,\rho(1-\eps/k)
 =\sqrt{M_0}\log(k/\eps)+O(\sqrt{M_0}),
 \qquad M_0=H\text{ or }L,
 \label{eq:formulation-route}
\end{equation}
for $0<\eps<k$, with the asymptotic form as $\eps/k\downarrow0$.
Sampling this path and the distance bounds give matching forward
$R$-chord counts $\Theta_R(\sqrt{M_0}\log(k/\eps))$, uniformly for $0<\eps\leq k/4$.
An approximate starting point and a fixed number of chords per round
are handled as in \eqref{eq:exposed-rounds}.
These statements concern the given barrier metrics; the block count
alone implies none of these primal lower bounds.
